% Auto-Compile LaTeX PDFs on Save
% Requirement: latexmk installed (via TeX Live or MacTeX).
%
% Start the watcher from the repo root:
%   ./watch-tex.sh
%
% It will build all .tex files once at startup, then watch for changes.
% Stop it with Ctrl+C. PDFs are generated alongside each .tex file, and
% intermediate build files are removed automatically after each build.

\documentclass[11pt]{article}
\usepackage[utf8]{inputenc}
\usepackage[T1]{fontenc}
\usepackage{amsmath,amssymb}
\usepackage{geometry}
\geometry{margin=1in}
\usepackage{enumitem}
\setlist[itemize]{leftmargin=*}
\setlist[itemize]{nosep}
\setlength{\parindent}{0pt}
\setlength{\parskip}{4pt}
\usepackage[hidelinks]{hyperref}

\newcommand{\RoadmapSection}[1]{%
  \section*{#1}%
  \addcontentsline{toc}{section}{#1}%
}

\newcommand{\RoadmapChapter}[1]{%
  \clearpage%
  \subsection*{#1}%
  \addcontentsline{toc}{subsection}{#1}%
}

\title{H2 Physics Learning Roadmap}
\date{\today}

\begin{document}
\maketitle
\clearpage
\tableofcontents
\clearpage

\RoadmapSection{SECTION 0: MATHEMATICAL FOUNDATION}

\RoadmapChapter{Chapter 0: Quantities and Analysis}

\begin{itemize}
  \item Units
  \begin{itemize}
    \item SI units: base units; derived units; unit prefixes
    \item Non-SI unit system; e.g.: Gaussian units, natural units
    \item Unit and estimation of quantities
    \item Units and dimensions
    \begin{itemize}
      \item Notation for quantity dimension
      \item Homogeneity of physical equations
      \item Dimensional analysis
    \end{itemize}
  \end{itemize}
  \item Measurements and Statistics
  \begin{itemize}
    \item Random and systematic errors of measurements
    \item Precision and accuracy of statistics
    \item Uncertainty calculations: absolute, fractional, percentage
  \end{itemize}
  \item Vector Algebra
  \begin{itemize}
    \item 2D vectors in coordinate systems
    \begin{itemize}
      \item Cartesian: $x$ and $y$ coordinates and orthogonal decomposition
      \item Polar: angular and radial coordinates
    \end{itemize}
    \item Scalars as 1D vectors
    \item Vector arithmetic
    \begin{itemize}
      \item Addition, multiplication (scalar, dot, and cross product)
      \item Mention: vector as an $n \times 1$ matrix, and the resulting algebra
    \end{itemize}
    \item Unit of vectors and vector products
    \item Mention: Complex numbers
    \begin{itemize}
      \item How are complex numbers similar to 2D vectors
      \item How they are different
    \end{itemize}
    \item Mention: fields
    \begin{itemize}
      \item Quantity as a function of spatial, or spatial and time, coordinates
      \item Superposition as vector addition
    \end{itemize}
  \end{itemize}
\end{itemize}

\RoadmapSection{SECTION I: MECHANICS}

\RoadmapChapter{Chapter 1: Kinematics I}

\begin{itemize}
  \item Position and its time derivatives
  \begin{itemize}
    \item Position, displacement, and distance
    \item Velocity and speed
    \item Acceleration and its magnitude
    \item Mention: jerk and higher order derivatives and why they are less often used
  \end{itemize}
  \item Graphs for 1D motions and analysis
  \begin{itemize}
    \item Gradients as the rate of change
    \item Areas as the accumulation
    \item Time dependent graphs: $x$--$t$, $v$--$t$, $a$--$t$
    \begin{itemize}
      \item Time derivative advances the curve
    \end{itemize}
    \item Briefly mention time independent graphs: $v$--$x$, $a$--$v$, $a$--$x$
    \begin{itemize}
      \item What advances the curve here?
    \end{itemize}
  \end{itemize}
  \item 1D motion with constant $\mathbf{a}$ (forces not mentioned)
  \begin{itemize}
    \item Equations of uniformly accelerated motion
    \begin{itemize}
      \item Derivation based on non-uniform motion
      \item Extra: Energy interpretation of the time independent equation
    \end{itemize}
    \item Example: Free fall (with constant gravitational acceleration)
  \end{itemize}
\end{itemize}

\RoadmapChapter{Chapter 2: Kinematics II}

\begin{itemize}
  \item 2D Motion overview
  \begin{itemize}
    \item Independent analysis along orthogonal directions
    \item With an interest over curvature: tangential and centripetal $\mathbf{a}$ by vector decomposition
  \end{itemize}
  \item 2D motion with constant $\mathbf{a}$ (forces not mentioned)
  \begin{itemize}
    \item Projectile motion
    \item Weight as force in a gravitational field ($W = mg$)
    \item Qualitative effects of air resistance, including terminal velocity
  \end{itemize}
  \item 2D motion with constant centripetal $|\mathbf{a}|$
  \begin{itemize}
    \item Angular displacement in radians
    \item Angular velocity $\omega$
    \item Relationship $v = r\omega$
    \item Common misconception: $\mathbf{a}$ is not constant for circular motion. $|\mathbf{a}|$ is.
    \item Centripetal acceleration ($a = \omega^2 r = v^2 / r$)
  \end{itemize}
\end{itemize}

\RoadmapChapter{Chapter 3: Dynamics I}

\begin{itemize}
  \item Mass as inertia (a qualitative mention)
  \item N1L: $\mathbf{F} = 0 \Rightarrow \mathbf{a} = 0$
  \item N2L: $\mathbf{F} = m\mathbf{a}$ (actually, $\mathbf{F}(t) = m(t)\mathbf{a}(t)$)
  \item N3L: $\mathbf{F}_{\text{action}} + \mathbf{F}_{\text{reaction}} = 0$
  \item Free-body diagrams and vector triangles
  \item Moment of a force, torque of a couple
  \item Conditions for equilibrium ($\mathbf{F}_{\text{net}} = 0$, $\tau_{\text{net}} = 0$) --- statics as special case
\end{itemize}

\RoadmapChapter{Chapter 4: Dynamics II}

\begin{itemize}
  \item Types of forces: gravitational, electric, magnetic, normal, friction, viscous, buoyant
  \item Centre of gravity
  \item Buoyancy
  \item Friction ($f_k = -\mu_k F_N \hat{v}$, $f_{s,\max} = -\mu_s F_N \hat{v}$)
  \item Drag (general model: polynomial with coefficients. more usually: linear / quadratic)
  \item Hooke's law ($\mathbf{F} = -k\Delta x$)
  \item Projectile and circular motion revisited
  \item Basic constrained motion --- kinematic relationships; inclined planes, pulleys, connected bodies
\end{itemize}

\RoadmapChapter{Chapter 5: Energy}

\begin{itemize}
  \item Energy stores and transfers, conservation of energy
  \item Work done by a force ($dW = \mathbf{F} \cdot d\mathbf{l} = |\mathbf{F}|\,|d\mathbf{l}|\cos\theta$)
  \item Kinetic energy --- derivation from work, $E_k = \tfrac12 m|\mathbf{v}|^2$
  \item Gravitational potential energy near Earth ($\Delta E_p = mg\Delta h$)
  \item Elastic potential energy --- area under force-extension graph
  \item Power ($P = E/t$, $P = Fv$)
  \item Energy transmission in constrained systems
  \item Efficiency
\end{itemize}

\RoadmapChapter{Chapter 6: Momentum}

\begin{itemize}
  \item Momentum --- $\mathbf{p} = m\mathbf{v}$
  \item Impulse --- $\mathbf{J} = \Delta \mathbf{p}$, area under force--time graph
  \item Conservation of momentum (Linear and Angular)
  \item Elastic and inelastic collisions in one dimension
  \begin{itemize}
    \item Relative speed of approach = relative speed of separation (perfectly elastic)
    \item Perfectly inelastic
    \item Kinetic energy changes in collisions
  \end{itemize}
  \item Extra: non-heads-on collisions
\end{itemize}

\RoadmapSection{SECTION II: OSCILLATIONS AND WAVES}

\RoadmapChapter{Chapter 7: Oscillations}

\begin{itemize}
  \item Free oscillations, experimental investigation
  \item Amplitude, period, frequency, angular frequency, phase, phase difference
  \item Defining equation of SHM ($a = -\omega^2 x$)
  \item General solution: ...
  \item Particular solutions: $x = x_0 \sin \omega t$, $v = v_0 \cos \omega t$, $v = \pm \omega\sqrt{x_0^2 - x^2}$
  \item Graphical relationships ($x$--$t$, $v$--$t$, $a$--$t$, $v$--$x$, $a$--$v$, $a$--$x$)
  \item Energy interchange in SHM
  \item Damped oscillations (light, critical, heavy), critical damping applications
  \item Forced oscillations and resonance --- amplitude response, sharpness, practical examples
\end{itemize}

\RoadmapChapter{Chapter 8: Waves I}

\begin{itemize}
  \item General wave equation
  \item Sinusoidal 1D wave as a function of $x$ and $t$
  \item Displacement, amplitude, period, frequency, wavelength, speed ($v = f\lambda$)
  \item Graphical representation --- transverse and longitudinal
  \item Energy transfer by waves, intensity ($I \propto \text{amplitude}^2$)
  \item Inverse square law for point sources
  \item Polarisation (transverse waves), Malus' law ($I = I_0 \cos^2 \theta$)
\end{itemize}

\RoadmapChapter{Chapter 9: Waves II}

\begin{itemize}
  \item Principle of superposition in 1D
  \item Standing waves for a general wave --- formation, nodes, antinodes (pressure vs displacement for sound)
  \item Examples: microwaves, strings, air columns, wavelength determination
  \item Qualitative understanding: Huygens' Principle
  \item Diffraction, interference, coherence, phase and path difference
  \item Two-source interference --- conditions, double-slit ($ax/D = \lambda$)
  \item Diffraction grating ($a\sin\theta = n\lambda$), wavelength determination
  \item Single-slit diffraction --- first minima ($b\sin\theta = \lambda$)
  \item Rayleigh criterion for resolution ($\theta \approx \lambda/b$)
\end{itemize}

\RoadmapSection{SECTION III: Gravity}

\RoadmapChapter{Chapter 10: Gravity}

\begin{itemize}
  \item Concept of field: quantity as a function of coordinates
  \item Newton's law of gravitation ($F = Gm_1m_2/r^2$) [change it to vector form]
  \item Gravitational field strength --- definition ($g = F/m$) and point mass formula ($g = GM/r^2$)
  \item Field lines --- uniform and radial
  \item Gravitational potential ($\phi = -GM/r$) and potential energy ($U = -GMm/r$)
  \item Field strength as negative potential gradient
  \item Escape velocity --- energy considerations
  \item Circular orbits --- equating gravitational force to centripetal force
  \item Geostationary satellites --- characteristics and applications
\end{itemize}

\RoadmapSection{SECTION IV: Electricity \& Magnetism}

\RoadmapChapter{Chapter 11: Electrostatics}

\begin{itemize}
  \item Electric field strength --- definition ($E = F/Q$) [Change everything to vector form]
  \item Coulomb's law ($F = (1/4\pi\varepsilon_0)(Q_1Q_2/r^2)$)
  \item Field strength for point charge ($E = (1/4\pi\varepsilon_0)(Q/r^2)$)
  \item Field lines and equipotentials
  \item Electric potential ($V = (1/4\pi\varepsilon_0)(Q/r)$) and potential energy ($U = (1/4\pi\varepsilon_0)(Q_1Q_2/r)$)
  \item Field strength as negative potential gradient
  \item Uniform fields between parallel plates ($E = V/d$), force on charge, motion of charged particles
  \item Capacitance ($C = Q/V$)
  \item Energy stored in a capacitor --- $U = \tfrac12 QV = \tfrac12 Q^2/C = \tfrac12 CV^2$, area under $V$--$Q$ graph
  \item Capacitors in series and parallel
\end{itemize}

\RoadmapChapter{Chapter 12: Circuit}

\begin{itemize}
  \item Electric current ($I = Q/t$), drift velocity ($I = nAvq$)
  \item Potential difference ($V = W/Q$)
  \item Electrical power ($P = VI = I^2R = V^2/R$)
  \item Distinguishing e.m.f. and p.d.
  \item Resistance ($R = V/I$), resistivity ($R = \rho l/A$)
  \item I--V characteristics --- ohmic resistor, diode, filament lamp, NTC thermistor
  \item Temperature dependence of resistivity --- metals (drift velocity), semiconductors (number density)
  \item Internal resistance --- effect on terminal p.d. and output power
  \item Resistors in series and parallel
  \item Potential divider circuits --- including thermistors and LDRs
  \item Charging and discharging a capacitor --- $Q = Q_0[1 - e^{-t/\tau}]$, $Q = Q_0e^{-t/\tau}$, time constant $\tau = RC$
\end{itemize}

\RoadmapChapter{Chapter 13: Magnetostatics}

\begin{itemize}
  \item Magnetic field from currents and permanent magnets
  \item Field lines --- long straight wire, flat circular coil, long solenoid
  \item Biot-Savart Law
  \item Magnetic flux density formulas: $B = \mu_0 I/(2\pi d)$, $B = \mu_0 NI/(2r)$, $B = \mu_0 nI$
  \item Force on current-carrying conductor ($F = BIl\sin\theta$), Fleming's left-hand rule
  \item Definition of magnetic flux density ($B = F/Il$ for $\perp$)
  \item Force between two current-carrying conductors
  \item Force on moving charge ($F = BQv\sin\theta$)
  \item Deflection of charged particles in electric and magnetic fields
  \item Velocity selection --- crossed fields
\end{itemize}

\RoadmapChapter{Chapter 14: Electromagnetic Induction}

\begin{itemize}
  \item Magnetic flux ($\Phi = BA$), flux linkage ($N\Phi$)
  \item Experimental evidence --- changing flux induces e.m.f.
  \item Faraday's law and Lenz's law --- direction opposes change
  \item Simple applications
  \item Alternating current --- period, frequency, peak value, r.m.s.
  \item Sinusoidal AC: $x = x_0 \sin \omega t$, $I_{\text{rms}} = I_0/\sqrt{2}$, $V_{\text{rms}} = V_0/\sqrt{2}$
  \item Mean power in resistive load --- half peak power for sine wave
\end{itemize}

\RoadmapChapter{Chapter 15: Electrical Appliances}

\begin{itemize}
  \item Half-wave rectification using a single diode
  \item Ideal transformer --- $N_s/N_p = V_s/V_p = I_p/I_s$, principle of operation
\end{itemize}

\RoadmapSection{SECTION V: THERMAL PHYSICS}

\RoadmapChapter{Chapter 16: Thermal Physics}

\begin{itemize}
  \item Thermodynamic temperature scale, absolute zero, conversion $^\circ$C to K
  \item Ideal gas equation: $pV = NkT$
  \item Avogadro constant, molar gas constant, Boltzmann constant
  \item Kinetic theory assumptions
  \item Derivation of $pV = \tfrac13 N m\langle c^2 \rangle$
  \item Mean translational KE $= \tfrac32 kT$
  \item Internal energy --- sum of microscopic KE and PE, determined by macroscopic state
  \item Temperature --- proportional to mean KE
  \item Thermal contact, thermal equilibrium, Zeroth law
  \item Work done by a gas ($W = p\Delta V$ at constant pressure)
  \item First law of thermodynamics ($\Delta U = Q + W$)
  \item Specific heat capacity and specific latent heat
\end{itemize}

\RoadmapSection{SECTION VI: MODERN PHYSICS}

\RoadmapChapter{Chapter 17: Quantum Physics}

\begin{itemize}
  \item Photoelectric effect --- threshold frequency, evidence for particle nature
  \item Photon energy ($E = hf$) and momentum ($p = h/\lambda$)
  \item Electron diffraction --- evidence for wave nature
  \item de Broglie wavelength ($\lambda = h/p$)
  \item Uncertainty principle
  \item Discrete electronic energy levels in atoms (e.g. hydrogen)
  \item Emission and absorption line spectra
  \item Photon transitions between energy levels
\end{itemize}

\RoadmapChapter{Chapter 18: Nuclear Physics}

\begin{itemize}
  \item Rutherford $\alpha$-scattering --- existence and small size of nucleus
  \item Nucleon number ($A$), proton number ($Z$), isotopes, notation $^A_Z X$
  \item Spontaneous and random nature of decay, fluctuations in count rate
  \item Background radiation --- origin and significance
  \item $\alpha$, $\beta$, $\gamma$ radiations --- nature and properties
  \item Activity ($A = \lambda N$), decay constant
  \item Exponential decay: $x = x_0 e^{-\lambda t}$ (activity, undecayed nuclei, count rate)
  \item Half-life: $t_{1/2} = \ln 2/\lambda$
  \item Applications and hazards of radioactivity (based on half-life, penetration, ionising ability)
  \item Nuclear equations and conservation laws (nucleon number, charge, mass-energy)
  \item Prediction of (anti)neutrino in $\beta$ decay (conservation of momentum and energy)
  \item Mass defect, binding energy, $E = mc^2$
  \item Binding energy per nucleon vs nucleon number --- graph and relevance to fusion and fission
\end{itemize}

\end{document}
